% Provenance (never printed): sources/book2/BRIEFS.md, g12-10. What a central bank aims at and controls; transmission and
% lags; rules and discretion (a Taylor-type rule as a description; time inconsistency; credibility and independence); the
% record (1979-83, the zero lower bound, 2021-23); one policy for many countries. Real numbers:
% figdata/grade-12/10-monetary-policy/ (Federal Reserve H.15 fed funds, H.4.1 total assets; BLS CPI-U; WDI GDP for the gap;
% JST unemployment 1979 and 1982; Brenland's invented accounts). Ledger: the Federal Reserve's policy-rules page and target-range
% table, the ECB key rates, the Bank of England's Bank Rate, Alesina-Summers 1993, and Book 3's rows on October 1979 and the
% 2008 and 2015 target ranges.
\chapter{Monetary Policy}\label{ch:g12:monetary-policy}

In June 1981 the average rate at which American banks lent to each other overnight reached 19.1\,\%. Fifteen months earlier, consumer prices had been rising at 14.8\,\% a
year. Twenty-seven years later, from December 2008, the same rate was held near zero for seven years. Both were decisions of the same institution, the Federal Reserve,
acting on the same instrument.
% figdata: policy_key (max_rate 19.1 at 1981.46 = June 1981; max_infl 14.76 at 1980.21 = March 1980); ledger: b1:fed-zlb-2008, b1:fed-liftoff-2015

This chapter asks what central banks try to achieve, what they control, how a change in their rate reaches the economy, and why a central bank that could choose freely
might do better by tying its own hands.

\begin{openingchart}{United States: effective federal funds rate (blue) and consumer price inflation over the previous twelve months (red), in per cent, monthly, January
1965 to October 2026.}{Board of Governors of the Federal Reserve System (H.15); U.S. Bureau of Labor Statistics (CPI-U, all items)}
\begin{tikzpicture}
\begin{axis}[omtime, width=0.9\linewidth, height=6.0cm, xmin=1965, xmax=2027, ymin=-3, ymax=20, ytick={0,5,...,20}, ylabel={per cent}, xtick={1970,1980,...,2020}]
\addplot[gray, thin, domain=1965:2027, samples=2] {0};
\addplot[omDef, thick] table[x=t, y=rate] {figdata/grade-12/10-monetary-policy/out/fedfunds.csv};
\addplot[omThm, thick] table[x=t, y=infl] {figdata/grade-12/10-monetary-policy/out/fedfunds.csv};
\node[font=\scriptsize, omDef, anchor=west] at (axis cs:1987,16) {federal funds rate};
\node[font=\scriptsize, omThm, anchor=west] at (axis cs:1993,-1.6) {inflation};
\end{axis}
\end{tikzpicture}
\end{openingchart}

\begin{recall}
From \cref{ch:g11:money-and-banking}: the central bank and its \omterm{def:g11:money-and-banking:policy-rate}{policy rate}. From \cref{ch:g10:price-indices}: the \omterm{def:g10:price-indices:real-interest-rate}{real interest rate}. From \cref{ch:g12:inflation-and-money}:
money and inflation in the long run. From \cref{ch:g12:aggregate-demand-and-supply}: \omterm{def:g12:aggregate-demand-and-supply:ad}{aggregate demand}, the \omterm{def:g12:aggregate-demand-and-supply:potential-output}{output gap}, and the return to the long run.
\end{recall}

\section{Aims and instruments}\label{sec:g12:monetary-policy:aims}

\begin{definition}[Inflation target]\label{def:g12:monetary-policy:inflation-target}
An \emph{inflation target}\index{inflation target} is a rate of inflation that a central bank announces it will aim at over time.
\end{definition}

Many central banks have one, often 2\,\% a year; some, like the Federal Reserve, are also asked to aim at high employment. The instrument is short: the central bank sets the
interest rate on overnight loans of \omterm{def:g11:money-and-banking:reserves}{central bank money}, and through it moves everything else.
% ledger: g12:boe-2026, g12:fed-rules

\begin{measure}[How a central bank sets an overnight rate]
% ledger: g12:ecb-2026
A central bank cannot order banks to lend to each other at its rate, but it can make other rates pointless. It pays a rate on the money banks leave with it overnight and
charges a higher rate on overnight loans it makes to them against collateral. No bank lends to another for less than it could earn at the central bank, and none borrows
for more than the central bank charges, so the market rate stays in the corridor between the two. From 16 September 2026 the European Central Bank paid 2.50\,\% on
deposits and charged 2.90\,\% on overnight loans, with its main lending rate at 2.65\,\%.
\end{measure}

\begin{dated}{2026}
% ledger: g12:fomc-2026, g12:ecb-2026, g12:boe-2026
\omterm{def:g11:money-and-banking:policy-rate}{Policy rates} in November 2026: Federal Reserve, a target range of 3.75--4.00\,\% since 17 September 2026; European Central Bank, 2.50\,\% on deposits since 16 September
2026; Bank of England, 4\,\% since 5 November 2026, with a target of 2\,\% inflation.
\end{dated}

\section{From the policy rate to the economy}\label{sec:g12:monetary-policy:transmission}

\begin{definition}[Monetary transmission]\label{def:g12:monetary-policy:transmission}
\emph{Monetary transmission}\index{monetary transmission} is the chain of effects by which a change in the \omterm{def:g11:money-and-banking:policy-rate}{policy rate} reaches spending, output and inflation.
\end{definition}

\begin{omfigure}
\begin{tikzpicture}[font=\scriptsize, node distance=0.45cm, box/.style={draw, rounded corners, fill=omDef!10, align=center, minimum height=0.9cm, text width=2.2cm}]
\node[box] (a) {policy rate rises};
\node[box, right=of a] (b) {market rates rise; asset prices fall};
\node[box, right=of b] (c) {borrowing, investment and consumption fall};
\node[box, right=of c, fill=omThm!10] (d) {output falls below potential};
\node[box, below=0.5cm of d, fill=omThm!10] (e) {inflation falls};
\node[box, below=0.5cm of b] (f) {currency rises; net exports fall};
\draw[->, thick] (a) -- (b); \draw[->, thick] (b) -- (c); \draw[->, thick] (c) -- (d); \draw[->, thick] (d) -- (e);
\draw[->, thick] (a.south) |- (f.west); \draw[->, thick] (f.east) -| (d.south west);
\end{tikzpicture}

{\small How a rise in the \omterm{def:g11:money-and-banking:policy-rate}{policy rate} reaches inflation: through the cost of borrowing, the prices of assets, and the exchange rate.}
\end{omfigure}

\begin{definition}[Policy lag]\label{def:g12:monetary-policy:lag}
A \emph{policy lag}\index{policy lag} is the delay between a change in the economy and its effects after a policy responds: the time to see the change in the data, to decide,
and for the decision to work through the chain.
\end{definition}

Because of the lags, a central bank must act on forecasts. A bank that waits until inflation is high to raise its rate, and until inflation has fallen to cut it, acts
too late both times.

\section{Rules and discretion}\label{sec:g12:monetary-policy:rules}

\begin{definition}[Policy rule and Taylor rule]\label{def:g12:monetary-policy:policy-rule}
A \emph{policy rule}\index{policy rule} is a formula that sets the \omterm{def:g11:money-and-banking:policy-rate}{policy rate} from data on the economy. A \emph{Taylor rule}\index{Taylor rule}%
\phantomsection\label{def:g12:monetary-policy:taylor-rule} sets the nominal rate at
\[ i = r^* + \pi + 0.5\,(\pi - \pi^*) + 0.5 \times \text{gap}, \]
where $r^*$ is a neutral real rate, $\pi$ inflation, $\pi^*$ the target and the gap is the \omterm{def:g12:aggregate-demand-and-supply:potential-output}{output gap} in per cent; the usual numbers are $r^* = 2$ and $\pi^* = 2$.
\end{definition}
% ledger: g12:fed-rules

The rule raises the nominal rate by more than any rise in inflation, so the real rate rises when inflation rises and falls when output is below potential.

\begin{omchart}{United States: the federal funds rate (blue) and the rate given by a Taylor rule (red), annual averages, 1965--2024, in per cent.}{Board of Governors of the
Federal Reserve System (H.15); U.S. Bureau of Labor Statistics; World Bank, World Development Indicators; rule computed in this book with consumer price inflation and the
\omterm{def:g12:aggregate-demand-and-supply:potential-output}{output gap} of \cref{ch:g12:aggregate-demand-and-supply}}
\begin{tikzpicture}
\begin{axis}[omtime, width=0.9\linewidth, height=5.4cm, xmin=1965, xmax=2024, ymin=-2, ymax=22, ytick={0,5,...,20}, ylabel={per cent}]
\addplot[gray, thin, domain=1965:2024, samples=2] {0};
\addplot[omDef, thick] table[x=year, y=actual] {figdata/grade-12/10-monetary-policy/out/taylor.csv};
\addplot[omThm, thick, dashed] table[x=year, y=rule] {figdata/grade-12/10-monetary-policy/out/taylor.csv};
\node[font=\scriptsize, omThm, anchor=south] at (axis cs:2014,13.5) {rule};
\node[font=\scriptsize, omDef, anchor=west] at (axis cs:1967,1) {actual};
\end{axis}
\end{tikzpicture}
\end{omchart}
% figdata: policy_key (corr_8707 0.69, mad_8707 1.5; rule2021 8.2 vs 0.08; rule2022 13.1 vs 1.68; rule2009 -1)

\begin{model}{A Taylor rule, as a description of policy}
\assumes The central bank reacts to inflation above its target and to the \omterm{def:g12:aggregate-demand-and-supply:potential-output}{output gap}, with fixed weights.
\predicts The path of the \omterm{def:g11:money-and-banking:policy-rate}{policy rate}, from data on inflation and output.
\tested As a description, not a test: in this book's version, from 1987 to 2007 the rule and the federal funds rate moved together (correlation 0.69), and the actual rate was on average 1.5 points away
  from the rule.
\breaks In 2009 the rule asked for a rate of $-1\,\%$, below zero; in 2021 it gave 8.2\,\% and in 2022 13.1\,\%, while the Federal Reserve kept its rate near zero until March
  2022. Its inputs, the neutral rate and \omterm{def:g12:aggregate-demand-and-supply:potential-output}{potential output}, are themselves estimates; policymakers consult such rules but do not follow them mechanically.
\end{model}

\begin{definition}[Time inconsistency]\label{def:g12:monetary-policy:time-inconsistency}
A policy is subject to \emph{time inconsistency}\index{time inconsistency} when a plan that is best before others act is no longer best for the policymaker once they have
acted.
\end{definition}

\begin{omfigure}
\centering
\small
\begin{tabular}{lcc}
\toprule
 & bank keeps inflation at 0\,\% & bank lets inflation reach 5\,\% \\
\midrule
wage setters expect 0\,\% & 0 & $+1$ \\
wage setters expect 5\,\% & $-3$ & $-1$ \\
\bottomrule
\end{tabular}

\medskip
{\small The central bank's \omterm{def:g11:oligopoly-and-games:game}{payoffs} in an invented \omterm{def:g11:oligopoly-and-games:game}{game}. Wage setters choose what to expect first, and want to be right; then the bank chooses inflation. A surprise
inflation raises output for a time ($+1$ net of the cost of inflation); a surprise of low inflation, after wages were set for 5\,\%, causes a slump ($-3$).}
\end{omfigure}

The bank would like everyone to expect 0\,\% and then keep inflation at 0\,\%, a \omterm{def:g11:oligopoly-and-games:game}{payoff} of 0. But once wages are set for 0\,\%, it gains by letting inflation reach 5\,\%.
Wage setters know this, expect 5\,\%, and the bank then lets inflation reach 5\,\% to avoid a slump: the outcome is high inflation and a \omterm{def:g11:oligopoly-and-games:game}{payoff} of $-1$, worse than the
plan.

\begin{definition}[Credibility and central bank independence]\label{def:g12:monetary-policy:credibility}
A central bank has \emph{credibility}\index{credibility} when people believe it will do what it announces. \emph{Central bank independence}\index{central bank independence}%
\phantomsection\label{def:g12:monetary-policy:independence} is the setting of policy by a central bank free of instructions from the government of the day.
\end{definition}

An independent bank with a clear target can be believed when it promises low inflation, because the government cannot order the surprise. Across sixteen rich countries,
central banks under political control were more likely to have high and variable inflation, and there was little sign that independence changed growth, unemployment or
\omterm{def:g10:price-indices:real-interest-rate}{real interest rates}. A correlation of this kind does not show by itself that independence lowers inflation: countries that dislike inflation may also choose independent
central banks.
% ledger: g12:as1993

\section{The record}\label{sec:g12:monetary-policy:record}

\begin{record}[The disinflation of 1979--83]
% ledger: b1:fed-targets; figdata: policy_key (pi1980 13.5, pi1983 3.2, actual1981 16.38, unemp1979 5.8, unemp1982 9.7)
On 6 October 1979 the Federal Reserve changed its procedure, targeting the growth of bank reserves to control the money supply, and let the overnight rate rise. The
federal funds rate averaged 16.4\,\% in 1981. Consumer price inflation fell from 13.5\,\% in 1980 to 3.2\,\% in 1983. The cost was a deep \omterm{def:g12:aggregate-demand-and-supply:recession}{recession}: unemployment rose from
5.8\,\% in 1979 to 9.7\,\% in 1982. In July 1982 the Federal Reserve's chairman told Congress that it was backing off its tight money targets. The expectations-augmented Phillips
curve of \cref{ch:g12:phillips-curve} predicts this cost when expectations are slow to fall: inflation came down only after a deep \omterm{def:g12:aggregate-demand-and-supply:recession}{recession}, and the
sacrifice ratio of that chapter's problem measures how deep.
\end{record}

\begin{definition}[Zero lower bound]\label{def:g12:monetary-policy:zlb}
The \emph{zero lower bound}\index{zero lower bound} is the floor, near zero, below which a central bank cannot push its \omterm{def:g11:money-and-banking:policy-rate}{policy rate} far, because people can hold cash, which
pays zero.
\end{definition}

On 16 December 2008 the Federal Reserve set a target range of 0 to $\tfrac14\,\%$, where it stayed until December 2015. With the rate at its floor, it used two other tools.

\begin{definition}[Quantitative easing and forward guidance]\label{def:g12:monetary-policy:qe}
\emph{Quantitative easing}\index{quantitative easing} is the purchase by a central bank of large amounts of long-term \omterm{def:g12:present-value:bond}{bonds} and other assets, paid for with new \omterm{def:g11:money-and-banking:reserves}{central bank
money}, to lower long-term interest rates. \emph{Forward guidance}\index{forward guidance}\phantomsection\label{def:g12:monetary-policy:forward-guidance} is a central bank's
announcement of how it expects to set its rate in the future.
\end{definition}

The Federal Reserve's total assets were 0.87 trillion dollars in August 2007, 4.5 trillion at the end of 2014 and 8.97 trillion at their peak in April 2022.
% ledger: b1:fed-zlb-2008, b1:fed-liftoff-2015; figdata: policy_key (assets_2007 0.87, assets_2014 4.5, assets_peak 8.97, 2022-04-13)

\begin{record}[2021--23]
% ledger: g12:fomc-2022; figdata: policy_key (pi2021 4.7, pi2022 8, pi2023 4.1, pi2024 2.9)
Consumer price inflation rose to 4.7\,\% in 2021 and 8.0\,\% in 2022. The Federal Reserve kept its target range at its floor until 17 March 2022, then raised it eleven
times, to 5.25--5.50\,\% on 27 July 2023. Inflation fell to 4.1\,\% in 2023 and 2.9\,\% in 2024. The rule of the chart had asked for higher rates a year earlier; whether
acting earlier would have cost less output, or whether inflation would have fallen anyway as supplies recovered, the record of one episode cannot say.
\end{record}

\section{One policy for many countries}\label{sec:g12:monetary-policy:union}

On 1 January 1999 the European Central Bank began setting one \omterm{def:g11:money-and-banking:policy-rate}{policy rate} for the countries of the euro area. One rate cannot fit every member at once: a country in a boom
with high inflation would want a higher rate than one in a slump. Within one country, too, regions differ, but they share a budget that moves money towards regions in
trouble and workers who move between them; how far a currency union can manage without these is a question for the chapters on trade and integration.
% ledger: g12:ecb-2026

\section{Exercises}

\begin{exercise}[$\star$]\label{exo:g12:monetary-policy:1}
From the opening chart, when did the federal funds rate reach its highest level, and roughly how high?
\end{exercise}

\begin{exercise}[$\star$]\label{exo:g12:monetary-policy:2}
From the opening chart, name the two periods since 2000 when the rate was near zero.
\end{exercise}

\begin{exercise}[$\star$]\label{exo:g12:monetary-policy:3}
With $r^* = 2$ and $\pi^* = 2$, find the Taylor-rule rate when inflation is 3\,\% and the \omterm{def:g12:aggregate-demand-and-supply:potential-output}{output gap} is $+1\,\%$.
\end{exercise}

\begin{exercise}[$\star$]\label{exo:g12:monetary-policy:4}
Find the Taylor-rule rate when inflation is 1\,\% and the \omterm{def:g12:aggregate-demand-and-supply:potential-output}{output gap} is $-2\,\%$.
\end{exercise}

\begin{exercise}[$\star$]\label{exo:g12:monetary-policy:5}
The \omterm{def:g11:money-and-banking:policy-rate}{policy rate} is 5.25\,\% and inflation 3\,\%. Find the real \omterm{def:g11:money-and-banking:policy-rate}{policy rate}, approximately.
\end{exercise}

\begin{exercise}[$\star\star$]\label{exo:g12:monetary-policy:6}
With a zero \omterm{def:g12:aggregate-demand-and-supply:potential-output}{output gap}, inflation is 2, 4 and then 6\,\% in three years. Find the rule's rate and the real rate it implies in each year.
\end{exercise}

\begin{exercise}[$\star\star$]\label{exo:g12:monetary-policy:7}
With inflation at 2\,\%, the \omterm{def:g12:aggregate-demand-and-supply:potential-output}{output gap} is $-4$, $-2$ and 0\,\% in three years. Find the rule's rate in each.
\end{exercise}

\begin{exercise}[$\star\star$]\label{exo:g12:monetary-policy:8}
Trace a cut in the \omterm{def:g11:money-and-banking:policy-rate}{policy rate} through the chain of transmission, link by link, to inflation.
\end{exercise}

\begin{exercise}[$\star\star$]\label{exo:g12:monetary-policy:9}
Inflation starts to rise in January. It shows in the data in March; the bank decides in June; a rate rise starts to lower inflation a year after it is made. When does
inflation start to fall because of the policy? What would a bank acting on forecasts do differently?
\end{exercise}

\begin{exercise}[$\star\star$]\label{exo:g12:monetary-policy:10}
From the chart of the rule, give the rule's rate and the actual rate in 2022, and the difference.
\end{exercise}

\begin{exercise}[$\star\star$]\label{exo:g12:monetary-policy:11}
In 2009 the rule gave $-1\,\%$. Why could the central bank not follow it, and what could it do instead?
\end{exercise}

\begin{exercise}[$\star\star\star$]\label{exo:g12:monetary-policy:12}
In 2021 the rule gave 8.2\,\% and the actual rate was close to zero. Does this show that the rule is wrong, that the central bank was wrong, or neither? Give two reasons a
bank might depart from the rule.
\end{exercise}

\begin{exercise}[$\star\star\star$]\label{exo:g12:monetary-policy:13}
In the \omterm{def:g11:oligopoly-and-games:game}{game} of this chapter, show that \emph{expect 0\,\%, keep 0\,\%} is not an equilibrium and that \emph{expect 5\,\%, let 5\,\%} is. How would a bank that could commit
in advance to 0\,\% do?
\end{exercise}

\begin{exercise}[$\star\star\star$]\label{exo:g12:monetary-policy:14}
State the cross-country finding on independence and inflation. Why does it not prove that making a central bank independent would lower inflation in a given country?
\end{exercise}

\begin{exercise}[$\star\star\star$]\label{exo:g12:monetary-policy:15}
Two regions of equal size share a currency. Region A has inflation of 4\,\% and an \omterm{def:g12:aggregate-demand-and-supply:potential-output}{output gap} of $+2\,\%$; region B has inflation of 0\,\% and a gap of $-3\,\%$. Find the rule's
rate for each region alone and for the union, using the average inflation and gap.
\end{exercise}

\section{Problem: Brenland's Central Bank}

\begin{problem}[{Weekend problem --- a rule against the record}]
\label{pb:g12:monetary-policy:1}
The table gives Brenland's \omterm{def:g11:money-and-banking:policy-rate}{policy rate}, consumer price inflation and \omterm{def:g12:aggregate-demand-and-supply:potential-output}{output gap} (real GDP against a log-linear trend fitted on 2001--2025, which grows by 1.52\,\% a year),
all in per cent. Brenland's central bank has an \omterm{def:g12:monetary-policy:inflation-target}{inflation target} of 2\,\%. Use the rule $i = 2 + \pi + 0.5(\pi - 2) + 0.5 \times \text{gap}$.

\medskip
\begin{center}
\small
\pgfplotstabletypeset[col sep=comma, columns={year,rate,infl,gap},
  columns/year/.style={column name=year, int detect, column type=l},
  columns/rate/.style={column name={\omterm{def:g11:money-and-banking:policy-rate}{policy rate}}, fixed, precision=2, zerofill},
  columns/infl/.style={column name=inflation, fixed, precision=1, zerofill},
  columns/gap/.style={column name={\omterm{def:g12:aggregate-demand-and-supply:potential-output}{output gap}}, fixed, precision=1, zerofill},
  every head row/.style={before row=\toprule, after row=\midrule}, every last row/.style={after row=\bottomrule}]
  {figdata/grade-12/10-monetary-policy/out/brenland.csv}
\end{center}
% figdata: brenland (rule and diff columns not printed)

\medskip
\noindent\textbf{Part I --- The record.}
\begin{enumerate}
\item In which year was inflation highest, and what was the \omterm{def:g11:money-and-banking:policy-rate}{policy rate} that year?
\item Find the real \omterm{def:g11:money-and-banking:policy-rate}{policy rate} in 2019 and in 2022.
\item In which year was the \omterm{def:g12:aggregate-demand-and-supply:potential-output}{output gap} most negative? What kind of event would explain it?
\item Was output above or below trend from 2016 to 2019?
\item Did the central bank raise its rate in 2022? By how much?
\end{enumerate}

\medskip
\noindent\textbf{Part II --- The rule.}
\begin{enumerate}[resume]
\item Find the rule's rate in 2015.
\item Find the rule's rate in 2020.
\item Find the rule's rate in 2022.
\item Find the gap between the rule's rate and the actual rate in 2022.
\item Find the rule's rate and the gap in 2023.
\end{enumerate}

\medskip
\noindent\textbf{Part III --- Reading the gap.}
\begin{enumerate}[resume]
\item The rule's rates for 2016--2019 are 5.65, 4.65, 5.15 and 4.95\,\%. Find the average gap between rule and actual in those years.
\item A rule that is too high in calm years may use the wrong neutral rate. With $r^* = 0$ instead of 2, what is the rule's rate in 2022?
\item What real rate does the rule ask for in 2022, and what was the actual real rate?
\item If \omterm{def:g12:aggregate-demand-and-supply:potential-output}{potential output} were higher than the trend, so that the 2022 gap was $-2\,\%$, what would the rule give?
\item Which inputs of the rule are measured, and which are estimated?
\end{enumerate}

\medskip
\noindent\textbf{Part IV --- Interpretation.}
\begin{enumerate}[resume]
\item Inflation fell to 5.6\,\% in 2023 and 3.8\,\% in 2024. Does this show that the bank was right not to follow the rule?
\item Give one reason, from the lags, why a bank might set its rate below the rule's when inflation is high.
\item Give one reason, from \omterm{def:g12:monetary-policy:credibility}{credibility}, why it might regret doing so.
\item Why might the rule's weights of 0.5 not suit Brenland?
\item State the result of the problem: the gap between the rule's rate and the actual \omterm{def:g11:money-and-banking:policy-rate}{policy rate} in 2022, and the inflation the rule would have reacted to.
\end{enumerate}
\end{problem}
